\hypertarget{ej5}{}
\noindent \textbf{Ejercicio 5.} Escribir una auxiliar que:

\begin{enumerate}[label=\alph*)]
    \item Sume 10 a un entero.
    \item ``Devuelva'' el mayor elemento de una tupla de 3 enteros.
    \item ``Devuelva'' el dígito menos significativo de un entero.
\end{enumerate}

\vspace{0.2cm}
{\color{gray}\hrule}
\vspace{0.4cm}

\textbf{Resolución:}

Recordemos que la sintaxis general para definir una función auxiliar en la materia es: \\
$\text{aux } \textit{nombre}(\textit{parametros}): \textit{tipo} = \textit{expresión}$

\begin{enumerate}[label=\textbf{\alph*)}]

    \item \textbf{Sume 10 a un entero} \\
    Simplemente tomamos un parámetro de tipo $\mathbb{Z}$ y le sumamos el valor constante.
    \[ \text{aux sumar10}(n: \mathbb{Z}): \mathbb{Z} = n + 10 \]

    \vspace{0.4cm}
    \item \textbf{Mayor elemento de una tupla de 3 enteros} \\
    Sea $t$ una tupla de tres elementos de tipo $\mathbb{Z} \times \mathbb{Z} \times \mathbb{Z}$. Usualmente accedemos a los elementos de una tupla mediante subíndices ($t_0, t_1, t_2$). Podemos utilizar la función \text{IfThenElse} para armar la lógica condicional:
    \begin{multline*}
    \text{aux maximoDeTres}(t: \mathbb{Z} \times \mathbb{Z} \times \mathbb{Z}): \mathbb{Z} = \\
    \text{IfThenElse}(t_0 \geq t_1 \wedge t_0 \geq t_2, \: t_0, \: \text{IfThenElse}(t_1 \geq t_2, \: t_1, \: t_2))
    \end{multline*}
    \textit{Nota: Otra forma 100\% válida (y más corta) si asumimos que ya existe una auxiliar básica \texttt{max(a,b)} es escribirlo como $\text{max}(t_0, \text{max}(t_1, t_2))$.}

    \vspace{0.4cm}
    \item \textbf{Dígito menos significativo de un entero} \\
    El dígito menos significativo de un número en base 10 se obtiene calculando el resto de su división entera por 10 (operación módulo). Para evitar resultados indeseados si el número llega a ser negativo, es una excelente práctica aplicarle el valor absoluto primero.
    \[ \text{aux digitoMenosSignificativo}(n: \mathbb{Z}): \mathbb{Z} = |n| \bmod 10 \]
    
\end{enumerate}

\vspace{0.5cm}
\hrule
\vspace{0.5cm}